% =============================================================================
% Assignments/CS301/06-linked-lists-singly-coding-assignment.tex
% Coding Assignment 6: P1 (Delete at Position) + P2 (Merge Sorted Lists)
%
% Student-facing problem statement. Instructor-only files live in
% Assignments/CS301/06-linked-lists-singly-coding/:
%   p1_reference_solution.c, p1_test_hidden.c   (NEVER published)
%   p2_reference_solution.c, p2_test_hidden.c   (NEVER published)
%   autograder.sh                                (NEVER published)
%   sll.h, sll_common.c, p1_starter.c, p1_test_visible.c,
%   p2_starter.c, p2_test_visible.c  (provided to students)
%
% Week 6 institute practical: singly-linked-list operations --
% syllabi/CS301.md Week 6 row ("instructor-supplied
% (singly-linked-list operations)"). P1 extends the lecture's
% deleteHead/deleteValue (Slides/CS301/06-linked-lists-singly.tex) to
% deletion by 0-based position, which the lecture never builds; P2 merges
% two sorted lists by rewiring, which the lecture never builds (and which
% Week 7 reuses for polynomials). Both use fresh lists of their own,
% distinct from the written assignment's 6->11->14, 6->11->9, 4->12->20,
% and 5->12->17->23 traces too.
%
% Compile with (Windows/MiKTeX, ';'-joined TEXINPUTS): cd Assignments/CS301 &&
%   TEXINPUTS="../../Preambles;" xelatex -interaction=nonstopmode 06-linked-lists-singly-coding-assignment.tex
%   (Windows/MiKTeX uses ; not : as the path separator)
% =============================================================================

\documentclass[11pt]{article}
\usepackage[margin=1in]{geometry}
% =============================================================================
% Preambles/header.tex — shared LaTeX/Beamer preamble for this project
%
% Use from a Beamer lecture by prepending:
%     \input{header}
% and compiling with TEXINPUTS=../Preambles:$TEXINPUTS (the /compile-latex
% skill does this automatically).
%
% PALETTE CONTRACT. Color names defined here MUST match the names in
% Quarto/theme-template.scss so Beamer and Quarto renderings use the same
% palette. The scripts/check-palette-sync.sh script enforces this.
%
% To customize: edit the HEX values below AND the matching $variable in
% Quarto/theme-template.scss, then run scripts/check-palette-sync.sh.
% =============================================================================

% -----------------------------------------------------------------------------
% Engine detection + encoding
%
% Our /compile-latex skill uses XeLaTeX, where inputenc is unnecessary and
% can produce encoding warnings. Only load inputenc under pdfTeX.
% -----------------------------------------------------------------------------
\usepackage{iftex}
\ifPDFTeX
  \usepackage[utf8]{inputenc}
\fi

\usepackage{amsmath, amssymb, amsthm}
\usepackage{booktabs}
\usepackage{graphicx}

% -----------------------------------------------------------------------------
% PALETTE — mirrors Quarto/theme-template.scss variable names.
% Load xcolor and define colors BEFORE anything that references them
% (notably hyperref, hypersetup, and the Beamer color assignments).
% -----------------------------------------------------------------------------
\usepackage{xcolor}

\definecolor{primary-blue}{HTML}{012169}      % headings, accents
\definecolor{primary-gold}{HTML}{B9975B}      % emphasis, borders
\definecolor{highlight-yellow}{HTML}{F2A900}  % markers, alerts
\definecolor{light-bg}{HTML}{E8EDF5}          % title / section backgrounds
\definecolor{jet}{HTML}{1A1A1A}               % body text

% Semantic colors (used in diagrams and annotations)
\definecolor{positive}{HTML}{15803D}          % good / observed / identified
\definecolor{negative}{HTML}{B91C1C}          % bad / problematic / bias
\definecolor{neutral}{HTML}{525252}           % reference / context

% Highlight accents (match .hi-* classes in the SCSS)
\definecolor{hi-slate}{HTML}{314F4F}
\definecolor{hi-green}{HTML}{2E8B57}
\definecolor{hi-red}{HTML}{C41E3A}

% Semantic aliases so skill templates and snippets can write
% \textcolor{accent}{...} without hard-coding "primary-blue".
\colorlet{accent}{primary-blue}
\colorlet{accent2}{primary-gold}

% -----------------------------------------------------------------------------
% Hyperref — loaded AFTER xcolor + palette so link colors resolve.
% -----------------------------------------------------------------------------
\usepackage{hyperref}
\hypersetup{colorlinks=true, linkcolor=primary-blue, urlcolor=primary-blue,
            citecolor=primary-blue, pdfborder={0 0 0}}

% -----------------------------------------------------------------------------
% Course code — set once per deck (after \input{header}, before \begin{document})
% via \coursecode{CS401}. Shown in the footer on every slide so a deck's course
% is identifiable without opening the title page. Empty by default.
% -----------------------------------------------------------------------------
\makeatletter
\newcommand{\coursecode}[1]{\gdef\@coursecodetext{#1}}
\gdef\@coursecodetext{}
\makeatother

% -----------------------------------------------------------------------------
% Beamer theme assignments (only apply under Beamer).
%
% We detect Beamer via \@ifclassloaded{beamer}, which is the documented,
% non-internal way. This must be wrapped in \makeatletter/\makeatother
% because \@ifclassloaded contains an @-catcode character.
% -----------------------------------------------------------------------------
\makeatletter
\@ifclassloaded{beamer}{%
  \usefonttheme{professionalfonts}
  \setbeamercolor{structure}{fg=primary-blue}
  \setbeamercolor{frametitle}{fg=primary-blue}
  \setbeamercolor{title}{fg=primary-blue}
  \setbeamercolor{subtitle}{fg=primary-gold}
  \setbeamercolor{author}{fg=primary-blue}
  \setbeamercolor{institute}{fg=neutral}
  \setbeamercolor{date}{fg=primary-gold}
  \setbeamercolor{itemize item}{fg=primary-blue}
  \setbeamercolor{itemize subitem}{fg=primary-gold}
  \setbeamercolor{alerted text}{fg=primary-gold}
  \setbeamercolor{block title}{fg=white, bg=primary-blue}
  \setbeamercolor{block body}{bg=light-bg}
  % Semantic block variants: block=definition/concept (blue), exampleblock=
  % worked example (green/positive), alertblock=key takeaway or caution (gold).
  % Keep to <=2 colored boxes per slide across ALL three variants (INV-7).
  \setbeamercolor{block title example}{fg=white, bg=positive}
  \setbeamercolor{block body example}{bg=light-bg}
  \setbeamercolor{block title alerted}{fg=white, bg=primary-gold}
  \setbeamercolor{block body alerted}{bg=light-bg}

  % Minimal footer: course code (left) + page X / N (right), no navigation clutter
  \setbeamertemplate{navigation symbols}{}
  \setbeamertemplate{footline}{%
    \color{neutral}\footnotesize\hspace{0.7em}\@coursecodetext\hfill
    \insertframenumber/\inserttotalframenumber\hspace{0.7em}\vspace{0.3em}}
}{}
\makeatother

% -----------------------------------------------------------------------------
% TikZ setup — libraries the snippet gallery uses
% -----------------------------------------------------------------------------
\usepackage{tikz}
\usetikzlibrary{arrows.meta, positioning, calc, decorations.pathreplacing,
                fit, shapes.geometric, backgrounds}

% Shared TikZ styles — snippets and hand-written diagrams can pick these up
% via `\begin{tikzpicture}[every node/.style={...}]` or by naming specific
% styles (e.g., `\node[dag-node] (X) at (0,0) {X};`).
\tikzset{
  % Node styles for causal DAGs and similar
  dag-node/.style = {
    draw=primary-blue, fill=light-bg, rounded corners=2pt,
    minimum width=1.2cm, minimum height=0.9cm, align=center, font=\small
  },
  decision-node/.style = {
    draw=primary-gold, fill=white, diamond, aspect=1.8,
    minimum width=1.8cm, minimum height=1.2cm, align=center, font=\small,
    inner sep=1pt
  },
  % Edge styles — observed vs counterfactual
  observed-edge/.style = {-{Latex[length=2.5mm]}, thick, draw=jet},
  counterfactual-edge/.style = {-{Latex[length=2.5mm]}, thick, draw=neutral, dashed},
  confound-edge/.style = {-{Latex[length=2.5mm]}, thick, draw=negative, dashed},
  % Data-point styles for plots
  observed-dot/.style = {fill=primary-blue, draw=primary-blue, circle, inner sep=1.2pt},
  counterfactual-dot/.style = {fill=white, draw=neutral, circle, inner sep=1.2pt}
}

% -----------------------------------------------------------------------------
% Convenience macros
% -----------------------------------------------------------------------------
\newcommand{\muted}[1]{\textcolor{neutral}{#1}}
\newcommand{\key}[1]{\textcolor{primary-gold}{\textbf{#1}}}
\newcommand{\good}[1]{\textcolor{positive}{#1}}
\newcommand{\bad}[1]{\textcolor{negative}{#1}}

% Standout frame macro: full-bleed dark background for section transitions.
% Beamer-only (uses \begin{frame}/\setbeamercolor/\usebeamerfont) -- do not
% call from an article-class file (e.g. Notes/<CODE>/*-notes.tex). \muted,
% \key, \good, \bad, and \coursecode above are plain xcolor/gdef and are
% safe to use from either class; verified when Notes/article-class support
% was added.
% Expand: \transitionslide{Your transition title}
\newcommand{\transitionslide}[1]{%
  \begingroup
  \setbeamercolor{background canvas}{bg=primary-blue}
  \begin{frame}[plain]
    \centering
    \vfill
    {\usebeamerfont{title}\color{white}\Large #1\par}
    \vfill
  \end{frame}
  \endgroup
}

% Section divider frame (Beamer-only), an alternative to \transitionslide with a
% darker two-line style: near-black (jet) background, a small white label line
% above a larger gold title, and a thin gold rule along the bottom edge.
% Expand: \sectiondivider{Section 2}{Pointers}
\newcommand{\sectiondivider}[2]{%
  \begingroup
  \setbeamercolor{background canvas}{bg=jet}
  \begin{frame}[plain]
    \vspace*{3.0cm}
    {\color{white}\ttfamily\large #1\par}
    \vspace{0.5cm}
    {\color{primary-gold}\LARGE #2\par}
    \begin{tikzpicture}[remember picture, overlay]
      \draw[primary-gold, line width=2.5pt]
        ($(current page.south west)+(0,0.1)$) -- ($(current page.south east)+(0,0.1)$);
    \end{tikzpicture}
  \end{frame}
  \endgroup
}

% -----------------------------------------------------------------------------
% Channel watermark — "Concepts That Click" (YouTube).
%
% Article-class files (CompetitiveExam Q/A sets, Notes, Assignments) get a
% light diagonal draft watermark on every page plus a centered footer naming
% the channel. Beamer decks get a subtle TikZ background watermark instead
% (the footline already carries the course code). Centralized here so every
% MCQ PDF is branded without per-file edits.
% -----------------------------------------------------------------------------
\makeatletter
\@ifclassloaded{article}{%
  \usepackage{draftwatermark}
  \SetWatermarkText{Concepts That Click}
  \SetWatermarkScale{1}
  \SetWatermarkFontSize{2cm}
  \SetWatermarkLightness{0.86}
  \SetWatermarkAngle{45}
  \usepackage{fancyhdr}
  \pagestyle{fancy}
  \fancyhf{}
  \fancyfoot[C]{\footnotesize\textcolor{neutral}{Concepts That Click~|~YouTube: \href{https://www.youtube.com/@concepts.thatclick}{@concepts.thatclick}}}
  \renewcommand{\headrulewidth}{0pt}
  \renewcommand{\footrulewidth}{0pt}
  \fancypagestyle{plain}{%
    \fancyhf{}%
    \fancyfoot[C]{\footnotesize\textcolor{neutral}{Concepts That Click~|~YouTube: \href{https://www.youtube.com/@concepts.thatclick}{@concepts.thatclick}}}%
    \renewcommand{\headrulewidth}{0pt}%
    \renewcommand{\footrulewidth}{0pt}%
  }
}{}
\@ifclassloaded{beamer}{%
  \setbeamertemplate{background}{%
    \begin{tikzpicture}[remember picture, overlay]
      \node[rotate=45, opacity=0.055, align=center]
        at (current page.center)
        {\Huge\textcolor{neutral}{Concepts That Click}};
    \end{tikzpicture}%
  }
}{}
\makeatother 
\coursecode{CS301}
\renewcommand{\thesection}{6.\arabic{section}}

\title{Coding Assignment 6: P1 -- Delete at Position, P2 -- Merge Sorted Lists}
\author{Auqib Hamid Lone\\[0.3em]
  \emph{Department of Computer Science \& Engineering}, \emph{GCET Ganderbal}}
\date{\today}

\begin{document}
\maketitle

\noindent\textbf{Due:} two weeks from issue. There is no automatic late penalty --- if you need more time, ask before the deadline and an extension will be granted (see the course policy in the syllabus).

\noindent This assignment autogrades \textbf{two} institute practicals on Week 6's singly linked list (\texttt{head} is the list handle, \texttt{p->next} is the successor link, \texttt{NULL} terminates the chain): \textbf{P1} (delete the node at a 0-based position) and \textbf{P2} (merge two sorted lists into one). The contract for both is given to you in \texttt{sll.h} --- \emph{do not modify it}. Helpers (\texttt{newNode}, \texttt{buildList}, \texttt{freeList}, \texttt{listToArray}, \texttt{listLength}) are implemented in \texttt{sll\_common.c} --- \emph{do not modify it either}. Each part is a separate program with its own starter file and hidden test suite.

\section{Why this problem}

The lecture built deletion by head and deletion by value, and traced reversal --- but never deletion \emph{by position} (drop the $k$-th entry of a log without knowing its value) and never \emph{merging} two ordered chains (combine two sorted runs without allocating anything new). Both are the same pointer surgery the lecture drilled --- save, unlink, free (P1); adopt-a-successor rewiring (P2) --- applied to the cases the lecture's own traces did not cover.

\section{Part 1 (P1): Delete at Position}

Implement in \texttt{p1.c} (start from \texttt{p1\_starter.c}):

\begin{verbatim}
int deleteAtPosition(struct node **head, int pos);
\end{verbatim}

\texttt{pos} is a 0-based index: \texttt{pos == 0} removes the head (the lecture's save/advance/free triple), \texttt{pos > 0} walks to the predecessor exactly as the lecture's \texttt{deleteValue} walk does, then unlinks and frees. Return $0$ on success. Return $-1$ and change \textbf{nothing} (no unlink, no free) when the list is empty, \texttt{pos} is negative, or \texttt{pos} is at or past the end of the list.

\subsection{Constraints}

\begin{itemize}
  \item The parameter must stay \texttt{struct node **head}: removing position $0$ reassigns the caller's \texttt{head}, exactly the lecture's double-pointer discipline.
  \item Unlink \emph{before} freeing --- the predecessor must point around the victim first (lecture's use-after-free rule).
  \item Time $O(n)$ worst case, $n$ = list length (the walk dominates); auxiliary space $O(1)$.
  \item Your code must compile with \texttt{gcc -Wall -Wextra} without errors or warnings, with no global/static mutable state.
\end{itemize}

\subsection{Worked examples (these are the visible tests)}

All on the list \texttt{head -> 7 -> 13 -> 19 -> 27 -> NULL}:

\begin{itemize}
  \item \texttt{deleteAtPosition(\&h, 1)} returns $0$; the list becomes \texttt{7 -> 19 -> 27 -> NULL}.
  \item \texttt{deleteAtPosition(\&h, 0)} returns $0$; the list becomes \texttt{13 -> 19 -> 27 -> NULL}.
  \item \texttt{deleteAtPosition(\&h, 3)} returns $0$; the list becomes \texttt{7 -> 13 -> 19 -> NULL}.
  \item \texttt{deleteAtPosition(\&h, 4)} returns $-1$; the list is unchanged.
  \item A singleton list \texttt{42 -> NULL} with \texttt{pos = 0} returns $0$ and leaves \texttt{h == NULL}.
\end{itemize}

\section{Part 2 (P2): Merge Two Sorted Lists}

Implement in \texttt{p2.c} (start from \texttt{p2\_starter.c}):

\begin{verbatim}
struct node *mergeSorted(struct node *a, struct node *b);
\end{verbatim}

Both inputs are ascending (either may be \texttt{NULL} for empty). \textbf{Reuse the existing nodes}: rewire \texttt{next} links only --- allocate nothing, free nothing, so every input node appears exactly once in the result. Return the head of the merged list (\texttt{NULL} only when both inputs are empty). On ties (\texttt{a->data == b->data}), take the node from \texttt{a} first. After the call, the old \texttt{a}/\texttt{b} heads are owned by the merged list --- do not use them.

\subsection{Constraints}

\begin{itemize}
  \item Time $O(n+m)$ worst case ($n, m$ = input lengths, each node visited once); auxiliary space $O(1)$ besides a few pointers.
  \item Your code must compile with \texttt{gcc -Wall -Wextra} without errors or warnings, with no global/static mutable state.
\end{itemize}

\subsection{Worked examples (these are the visible tests)}

\begin{itemize}
  \item \texttt{mergeSorted([2, 8, 21], [5, 15, 33])} returns\\ \texttt{2 -> 5 -> 8 -> 15 -> 21 -> 33 -> NULL}.
  \item \texttt{mergeSorted([3, 11, 29], [])} returns \texttt{3 -> 11 -> 29 -> NULL}.
  \item \texttt{mergeSorted([], [])} returns \texttt{NULL}.
  \item \texttt{mergeSorted([4, 9, 18], [9, 18, 26])} returns\\ \texttt{4 -> 9 -> 9 -> 18 -> 18 -> 26 -> NULL}.
\end{itemize}

\section{What to submit for each part}

\begin{enumerate}
  \item \textbf{P1:} \texttt{p1.c} (your completed implementation) and the unmodified \texttt{sll.h}, \texttt{sll\_common.c}.
  \item \textbf{P2:} \texttt{p2.c} (your completed implementation) and the unmodified \texttt{sll.h}, \texttt{sll\_common.c}.
\end{enumerate}

Sanity-check locally before submitting:
\begin{verbatim}
gcc -Wall -Wextra -o p1_test p1_test_visible.c p1.c sll_common.c && ./p1_test
gcc -Wall -Wextra -o p2_test p2_test_visible.c p2.c sll_common.c && ./p2_test
\end{verbatim}
All worked examples above should pass. The autograder then compiles each part against its own \emph{hidden} test suite. Each hidden suite covers, at minimum:
\begin{itemize}
  \item the worked examples above, re-checked on fresh data;
  \item boundary sizes (empty/singleton/negative-position for P1; empty/singleton/tie for P2);
  \item structure-specific adversarial cases --- for P1: head/tail/out-of-range/huge-position, draining a list to empty, all-duplicate lists, a 100-element list; for P2: one side entirely smaller, unequal lengths, negative values, duplicates within and across lists, a 500-node merge, sortedness plus node conservation on every merge;
  \item and a batch of randomly generated inputs checked against an independent oracle.
\end{itemize}

\end{document}
